\section{CONCLUSÃO}

A investigação termodinâmica e estrutural conduzida neste trabalho evidenciou a complexa rede de forças que governa a estabilidade e a dinâmica de fases das membranas lipídicas. Utilizando simulações de Dinâmica Molecular com o campo de força \textit{coarse-grained} Martini 3, foi possível não apenas capturar as propriedades macroscópicas esperadas, mas também dissecar quantitativamente os domínios de energia eletrostática (Coulomb) e de dispersão (Lennard-Jones) subjacentes a esses fenômenos.

A análise sob diferentes gradientes de hidratação revelou uma marcante assimetria cinética entre fosfolipídios homólogos. Demonstrou-se que a matriz de DSPC ($C_{18}$), ao contrário do DPPC ($C_{16}$), sucumbe a uma armadilha metaestável sob alta hidratação, manifestando-se como um líquido super-resfriado. Esta retenção foi termodinamicamente associada à barreira entálpica de dessolvatação interfacial: a alta inércia estérica das cadeias maiores exige uma energia de ativação cooperativa que suplanta as flutuações térmicas disponíveis, impedindo a cristalização espontânea para a fase gel.

A modulação imposta pela inserção de colesterol elucidou o funcionamento de um mecanismo de compensação isoperibólica de extrema elegância biofísica. A drástica condensação lateral induzida pelas altas concentrações de esterol força a expulsão da água interfacial, o que acarreta um severo déficit nas interações de Coulomb e na hidratação primária. Contudo, demonstrou-se que a energia global de dispersão do sistema se mantém estável, pois o custo da dessolvatação é integralmente neutralizado pelo ganho de contatos atrativos de van der Waals decorrentes do empacotamento linear (\textit{all-trans}) das caudas hidrofóbicas contra a estrutura rígida do esterol.

Adicionalmente, a inserção de íons no sistema evidenciou o papel da coordenação do sódio ($Na^+$) na estabilização adicional da interface polar. Contudo, em misturas saturadas com colesterol (40\% em mol), a rigidez mecânica imposta pela fase líquida-ordenada ($L_o$) torna-se dominante, anulando a cooperatividade intrínseca da matriz lipídica e sobrepondo-se até mesmo aos fortes efeitos de ancoragem iônica.

Os resultados aqui consolidados transcendem o entendimento fundamental da física de membranas. Ao comprovar a sensibilidade topológica e a precisão do modelo Martini 3 no balanço de interações não-covalentes, este trabalho pavimenta um alicerce para a biofísica translacional. A elucidação da dinâmica interfacial e do empacotamento lipídico estabelece diretrizes termodinâmicas rigorosas para o \textit{design} de lipossomas na vetorização de fármacos e para a prospecção \textit{in silico} de agentes antivirais, consolidando a Dinâmica Molecular como uma ferramenta preditiva indispensável e de altíssimo rigor para o avanço da física médica.
